% Source ownership:
% - Itxaso

% ============================================================
\section{MPI Parallelization: Observable Post-Processing}
\label{sec:parallel_observables}
% ============================================================

\subsection{Design and Implementation}

Post-processing of the trajectory data has been parallelized in
\texttt{main\_parallel\_observables.f90}. Rather than decomposing a single
trajectory into contiguous frame blocks, the implementation operates at the
file level: all trajectory files are first enumerated in a manifest, rank 0
reads this manifest, and the complete file list is broadcast to all MPI ranks.
Each rank then processes a round-robin subset of trajectory files,
independently computing the radius of gyration \(R_g\), the end-to-end
distance \(R_{ee}\), and the torsional statistics for the configurations
assigned to it.

This design was chosen because the available data are naturally distributed
across multiple trajectory files, corresponding to different configuration
types and simulation phases (\texttt{equil} and \texttt{prod}). A file-level
distribution avoids repeated file seeks inside large XYZ trajectories, keeps
the implementation simple, and preserves a fully embarrassingly parallel local
workload. At the same time, it remains fully compatible with collective MPI
communication for global aggregation.

A first collective step is performed through \texttt{MPI\_Bcast}, which is used
to distribute the manifest and the maximum number of torsional degrees of
freedom to all ranks. After each rank has completed its local file subset, a
mid-run checkpoint is evaluated by means of \texttt{MPI\_Allreduce}. In this
step, the partial frame counts and partial \(R_g\) sums are synchronised across
all processes, so that the global partial mean can be computed before the final
aggregation stage. This checkpoint was introduced as a lightweight example of
active communication during the computation phase, beyond the final collection
of results.

\begin{lstlisting}[caption={Round-robin file distribution and MPI checkpoint in \texttt{main\_parallel\_observables.f90}.}]
! Each rank processes a round-robin subset of trajectory files
do ifile = rank + 1, nfiles, num_procs
  call process_trajectory(trim(files(ifile)), max_phi, &
       local_file_count, local_frame_count, &
       local_sum_rg, local_sum_rg2, local_sum_ree, local_sum_ree2, &
       local_nphi, local_sum_phi, local_sum_phi2)
enddo

! Mid-run checkpoint: all ranks obtain the global partial mean of Rg
call MPI_Allreduce(local_frame_count, ckpt_frame_count, n_conf*n_phase, &
                   MPI_INTEGER, MPI_SUM, MPI_COMM_WORLD, ierr)
call MPI_Allreduce(local_sum_rg, ckpt_sum_rg, n_conf*n_phase, &
                   MPI_DOUBLE_PRECISION, MPI_SUM, MPI_COMM_WORLD, ierr)

! Final aggregation on rank 0
call MPI_Reduce(local_sum_rg, global_sum_rg, n_conf*n_phase, &
                MPI_DOUBLE_PRECISION, MPI_SUM, 0, MPI_COMM_WORLD, ierr)
\end{lstlisting}

The final global statistics are obtained on rank 0 through a series of
\texttt{MPI\_Reduce} operations, which combine the local sums, squared sums,
frame counts, and torsional accumulators into global averages and standard
deviations. The corresponding summary files
(\texttt{summary\_global\_np1.dat}, \texttt{summary\_global\_np2.dat},
\texttt{summary\_global\_np4.dat}, and \texttt{summary\_global\_np8.dat}) show
that the numerical results are identical for all tested values of \(n_p\),
confirming that the parallelization preserves the observable averages exactly.

The benchmark driver is invoked from the cluster submission script
\texttt{3.run\_parallel\_observables\_np.sh}, which executes the same
post-processing workflow for \(n_p \in \{1,2,4,8\}\), records both MPI and
shell wall-clock times, and writes per-process summary files for subsequent
comparison.

% ============================================================